\section{含有两个和三个味的 $\SU(2)$ 理论}\label{sec:nf=2,3}
\subsection{一般讨论}
本节及下一节考虑含有 $N_f$ 个味的 $\SU(2)$ 规范理论，其中 $N_f=2,3,4$。用 $\cN{=}1$ 手征多重态表示时，场为 $(Q_i,\tilde Q^i)$，其中 $i=1,\ldots,N_f$，超势为
\begin{equation}
\sum_i\left(Q_i\Phi\tilde Q^i+\mu_iQ_i\tilde Q^i\right)
\end{equation}
其中 $\mu_i$ 是裸质量项。当所有 $\mu_i$ 相同时，存在作用于 $Q_i$ 和 $\tilde Q^i$ 的指标 $i$ 上的 $\U(N_f)$ 对称性。在 $\Phi=\diag(a,-a)$ 的库仑支上，超多重态的物理质量为
\begin{equation}
|{}\pm a\pm\mu_i|.
\end{equation}

当 $\mu_i=0$ 时，可以将 $Q_i$ 和 $\tilde Q^i$ 合并为
\begin{equation}
(q_I^a)_{I=1,2,\ldots,2N_f}=(Q_1^a,\ldots,Q_{N_f}^a,\epsilon^{ab}\tilde Q^1_b,\ldots,\epsilon^{ab}\tilde Q^{N_f}_b)
\end{equation}
并具有 $\SO(2N_f)$ 对称性。在这一记号下，超势为
\begin{equation}
\propto\eta^{IJ}q_I^a\Phi_{ab}q_J^b,\qquad
\text{其中}\ \eta=\left(
\begin{array}{c|c}
0&\mathbf{1}_{N_f}\\\hline
\mathbf{1}_{N_f}&0
\end{array}\right).
\end{equation}
由于 $\eta^{IJ}$ 是对称矩阵，作用于指标 $I,J$ 的味对称性就是 $\SO(2N_f)$。等价地说，我们有 $2N_f$ 个处于 $\SU(2)$ 二重态表示中的半超多重态。

在经典层面，引入奇数个 $\SU(2)$ 二重态半超多重态没有问题，此时具有 $\SO(\text{odd})$ 味对称性。然而，这样的理论会有奇数个二重态 Weyl 费米子，因而在量子层面受到 Witten 整体反常的困扰，正如第 \ref{sec:anomaly} 小节所回顾的。因此，对于 $\SU(2)$ 规范群，只能考虑偶数个二重态半超多重态，或等价地，整数个二重态完整超多重态。

这一理论在紫外区域 $|a|\gg|\mu_i|$ 中的单圈跑动为
\begin{equation}
\tau(a)=2\tau_{UV}-\frac{2(4-N_f)}{2\pi i}\log\frac{a}{\Lambda_{UV}}+\cdots\label{su2running}
\end{equation}
当 $N_f\neq4$ 时，还可以进一步重写为
\begin{equation}
=-\frac{2(4-N_f)}{2\pi i}\log\frac{a}{\Lambda}\label{general_1_loop}
\quad\text{其中}\quad
\Lambda^{4-N_f}=\Lambda_{UV}^{4-N_f}e^{2\pi i\tau_{UV}}.
\end{equation}

我们在第 \ref{sec:su2nfpre} 小节猜出了这些理论的曲线形式。对于 $N_f=2,3,4$，结果分别由式 \eqref{nf_2_curve_pre}、\eqref{nf_3_curve_pre}、\eqref{nf_4_curve_pre} 给出。本节及下一节的目的，是进行各种检验，确认它们确实重现预期性质，并利用它们研究强耦合动力学。本节处理 $N_f=2$ 和 $N_f=3$ 的情形。$N_f=4$ 则打开了一个全新的领域，第 \ref{sec:gaiotto} 节将专门讨论它。

\subsection{$N_f=2$：曲线与单值变换}\index{SU2TwoFlavors@含两个味的 $\SU(2)$ 理论|hyperpage}
从含有 $N_f=2$ 个味的 $\SU(2)$ 理论开始。\SeibergWitten\ 曲线已在式 \eqref{nf_2_curve_pre} 中猜出，在此重写为
\begin{equation}
\Sigma:\qquad\frac{2\Lambda(x-\mu_1)}z+2\Lambda(x-\mu_2)z=x^2-u\label{nf_2_curve}
\end{equation}
配以 Seiberg-Witten 微分 $\lambda=xdz/z$。\Gaiotto 曲线 $C$ 仍然只是一个 $S^2$，如图 \ref{fig:nf_2_curve} 所示。
\begin{figure}[h]
\[
\inc{su2nf2-curve}
\]
\caption{$N_f=2$ 理论的曲线。\label{fig:nf_2_curve}}
\end{figure}

当 $|u|\gg|\Lambda^2|,|\mu_i|^2$ 时，很容易估算线积分。首先将 $A$ 闭链放在 $|z|=1$ 处。利用这一带的 $x\simeq\sqrt{u}$，得到
\begin{equation}
a=\frac{1}{2\pi i}\oint_Ax\frac{dz}z\simeq\sqrt{u}.\label{nf_2_a}
\end{equation}
$x(z)$ 在曲线 $C$ 上的分支点位置也很容易估算：两个位于 $z\simeq\sqrt{u}/\Lambda$ 附近，另两个位于 $z\simeq\Lambda/\sqrt{u}$ 附近。因而
\begin{equation}
a_D=\frac{1}{2\pi i}\oint_Bx\frac{dz}z\simeq\frac{2\cdot2}{2\pi i}\int^1_{\sqrt{u}/\Lambda}a\frac{dz}z
\simeq-\frac{4}{2\pi i}a\log\frac{a}{\Lambda}.\label{nf_2_ad}
\end{equation}
由此计算 $\tau(a)=\partial a_D/\partial a$，就重现了单圈跑动 \eqref{general_1_loop}。

接下来研究 $u$ 平面上的奇点结构。当 $\mu_{1,2}\gg\Lambda$ 时，能标 $\mu_{1,2}$ 附近的规范耦合相当小。因此预期，当
\begin{equation}
u\simeq\mu_i^2\quad\text{其中 $i=1,2$}
\end{equation}
时，$(Q_i,\tilde Q^i)$ 的一个分量变得非常轻，从而产生一个奇点。在能标 $\mu_{1,2}$ 以下，理论有效地等价于纯 $\SU(2)$ 理论，应当有两个奇点，分别对应单极子或双荷子变得非常轻。因此，总共预期 $u$ 平面上有四个奇点，见图 \ref{fig:nf_2_u_plane_generic}。
\begin{figure}[h]
\[
\inc{su2nf2-uplane-generic}
\]
\caption{$N_f=2$ 理论的 $u$ 平面。\label{fig:nf_2_u_plane_generic}}
\end{figure}

从曲线 \eqref{nf_2_curve} 出发，研究其判别式即可验证这一结构，这留给读者作为练习。这里详细研究无质量情形 $\mu_1=\mu_2=0$。此时的 \SeibergWitten\ 曲线就是
\begin{equation}
x^2-2\Lambda(z+\frac1z)x-u=0.
\end{equation}
可以看到，当 $u=0$ 或 $u=-4\Lambda^2$ 时，$x(z)$ 的四个分支点两两相遇，如图 \ref{fig:collision_nf_2} 所示。
\begin{figure}[h]
\[
\inc{su2nf2-curve-degenerate}
\]
\caption{当 $u=0$ 或 $u=-4\Lambda^2$ 时，$N_f=2$ 理论的曲线发生退化。\label{fig:collision_nf_2}}
\end{figure}

具体而言，当 $u=0$ 时，它们在 $z=\pm i$ 处相遇；当 $u=-\Lambda^2$ 时，它们在 $z=\pm1$ 处相遇。$u$ 平面上没有其他奇点，因此，当 $\mu_1=\mu_2=0$ 时，$u$ 平面具有图 \ref{fig:u_plane_nf_2_massless} 所示的结构。

% [Tutor Note: 原文在 p.80 同一页内即已前后不一致：叙述"four branch points … meet in pairs when $u=0$ or $u=-4\Lambda^2$"以及图 \ref{fig:collision_nf_2} 的图注均作 $u=-4\Lambda^2$，紧接着的 "Explicitly, …" 一句却改作 $u=-\Lambda^2$，其后（pp.81–83）一律沿用 $u=-\Lambda^2$。由无质量曲线 $x^2-2\Lambda(z+1/z)x-u=0$ 可直接验算：$u=0$ 时分支点在 $z=\pm i$ 相遇、$u=-4\Lambda^2$ 时在 $z=\pm1$ 相遇，而 $u=-\Lambda^2$ 时并不相撞，故 $u=-4\Lambda^2$ 为正确值（v3 已统一为 $-4\Lambda^2$）。本译本按底本 v2 照排。]
\begin{figure}[h]
\[
\inc{su2nf2-uplane-massless}
\]
\caption{无质量 $N_f=2$ 理论的 $u$ 平面。\label{fig:u_plane_nf_2_massless}}
\end{figure}
在 $u=0$ 和 $u=-\Lambda^2$ 这两处，分别都有两对 $x(z)$ 的分支点碰撞。这意味着，$u=0$ 和 $u=-\Lambda^2$ 每一处都应当视为 $u$ 平面上的两个奇点。图 \ref{fig:u_plane_nf_2_massless} 通过在 $u=0,-\Lambda^2$ 各放置两个几乎重叠的圆点，表示了这种情形。总共存在四个奇点，与上面对 $\mu_{1,2}\gg\Lambda$ 的结论相符。我们如图 \ref{fig:u_plane_nf_2_massless} 所示，记下绕各条闭路径的单值变换。

从式 \eqref{nf_2_a}、\eqref{nf_2_ad} 中 $a$、$a_D$ 的显式形式，可以得到无穷远处的单值变换
\begin{equation}
M_\infty=\begin{pmatrix}-1&2\\0&-1\end{pmatrix}.
\end{equation}
为求出绕 $u=0$ 的单值变换 $M_+$，可以取很小的 $\epsilon$，沿路径 $u=\epsilon e^{i\theta}$ 从 $\theta=0$ 缓慢改变至 $\theta=2\pi$，并跟踪分支点的运动。
\begin{figure}[h]
\[
\inc{su2nf2-cycles-monodromy}
\]
\caption{$N_f=2$ 理论中单值变换对闭链的作用。\label{fig:nf_2_mon}}
\end{figure}
分支点对如图 \ref{fig:nf_2_mon} 所示交换位置。可以看到，$B$ 闭链保持不变，而 $A$ 被映为 $A-2B$，从而产生
\begin{equation}
M_+=\begin{pmatrix}1&0\\-2&1\end{pmatrix}=ST^2S^{-1}.\label{nf_2_m+}
\end{equation}
我们看到
\begin{equation}
M_\infty=M_+M_-
\end{equation}
其中
\begin{equation}
M_-=TM_+T^{-1}.
\end{equation}

在继续之前，利用曲线的另一种描述得到相同的 $u$ 平面结构，也很有启发。该曲线由式 \eqref{nf_2_curve_pre2} 给出。在无质量情形下，它变为
\begin{equation}
\frac{x^2}z+\Lambda^2z=x^2-u.
\end{equation}
与此前一样，分支点在 $u=0$ 或 $u=-\Lambda^2$ 时发生碰撞，不过在\Gaiotto 曲线上看起来颇为不同，如图 \ref{fig:nf_2_second} 所示。
\begin{figure}[h]
\[
\inc{su2nf2-curveprime-degenerationX}
\]
\caption{第二种描述：当 $u=0$ 或 $u=-\Lambda^2$ 时，$N_f=2$ 理论的曲线发生退化。\label{fig:nf_2_second}}
\end{figure}

注意，不论 $u$ 取何值，$\lambda$ 都在 $z=0$ 和 $z=1$ 处发散。\Gaiotto 曲线上 $x(z)$ 的一个分支点随着 $u$ 改变而移动；当 $u=0$ 或 $u=-\Lambda^2$ 时，这个分支点分别撞上 $z=0$ 或 $z=1$。从后一种观点重现单值变换 $M_\pm$，留给读者完成。

\subsection{$N_f=2$：离散 R 对称性}\index{DiscreteRSymmetry@离散 R 对称性|hyperpage}\index{FlavorParity@味宇称|hyperpage}\index{CompositeOperator@复合算符|hyperpage}
现在研究离散 R 对称性。将连续 R 对称性下的荷赋值为
\begin{equation}
\begin{array}{r|ccc}
R=0&&A\\1&\lambda&&\lambda\\2&&\Phi
\end{array},\qquad
\begin{array}{r|ccc}
R=-1&&\psi_I\\0&q_I&
\end{array}.
\end{equation}
旋转
\begin{equation}
\lambda\to e^{i\varphi}\lambda,\qquad
\psi_I\to e^{-i\varphi}\psi_I\label{nf_2_shiftx}
\end{equation}
具有反常，但可以通过
\begin{equation}
\theta_{UV}\to\theta_{UV}+4\varphi.\label{nf_2_shift}
\end{equation}
予以补偿。等价地，动力学能标 $\Lambda$ 按照
\begin{equation}
\Lambda^2\to e^{4i\varphi}\Lambda^2.
\end{equation}
变换。因此，$\varphi=\pi/2$ 是一个真正的对称性，其作用为
\begin{equation}
\theta_{UV}\to\theta_{UV}+2\pi,\quad
\Phi\to-\Phi,\quad
u'\to u'
\end{equation}
其中 $u'=\vev{\tr\Phi^2/2}$。为何这里给符号 $u$ 加上撇号，稍后就会解释。遗憾的是，这一对称性在 $u'$ 平面上作用平凡，因此并不能告诉我们有关该平面结构的太多信息。

可以进行一种稍微更微妙的操作。考虑对超多重态的如下作用：
\begin{align}
(q_{I=1},q_2,q_3,q_4)&\mapsto(-q_{I=1},q_2,q_3,q_4),\\
(\psi_{I=1},\psi_2,\psi_3,\psi_4)&\mapsto(-\psi_{I=1},\psi_2,\psi_3,\psi_4).\label{flavor_parity}
\end{align}
迄今我们一直说味对称性是 $\SO(2N_f)=\SO(4)$。这一操作是味宇称作用，
\begin{equation}
\diag(-1,+1,+1,+1)\in\OO(4)\supset\SO(4).
\end{equation}
回忆一下，在 $\SU(2)$ 的 $k$ 瞬子背景中，$\psi_{I=1}$ 的零模数恰好为 $k$。因而，操作 \eqref{flavor_parity} 使路径积分测度乘以 $(-1)^k$。这意味着经典味对称性 $\mathrm{O}(4)$ 的宇称部分具有反常。不过，路径积分的被积函数中有一项 $e^{i\theta k}$，所以可以通过平移 $\theta\to\theta+\pi$ 来补偿它。

因此，可以将 $\varphi=\pi/4$ 的相位旋转 \eqref{nf_2_shiftx}、\eqref{nf_2_shift} 与味宇称 \eqref{flavor_parity} 组合起来，得到一个真正未破缺的对称性。总结起来，它是两个作用的组合：第一个是
\begin{equation}
\theta\to\theta+\pi,\quad\Phi\to i\Phi,\quad u'\to-u'
\end{equation}
第二个是
\begin{equation}
\theta+\pi\to\theta+2\pi,\quad q_{I=1}\to-q_{I=1},\quad\psi_{I=1}\to-\psi_{I=1}.
\end{equation}
总的来说，这是一个 $\bZ_4$ 对称性，它通过 $\bZ_2$ 作用于 $u$ 平面。

乍看之下，这似乎与图 \ref{fig:u_plane_nf_2_massless} 中得到的 $u$ 平面结构矛盾：两个奇点位于 $u=0$ 和 $u=-\Lambda^2$。解决办法是令
\begin{equation}
u=u'-\frac{\Lambda^2}{2}.\label{ushift}
\end{equation}

% [Tutor Note: 若采用上注的正确奇点 $u=0$、$u=-4\Lambda^2$，并要求离散 R 对称性在 $u'$ 上作用为 $u'\to-u'$（见前面 $\varphi=\pi/2$ 的作用），则两个奇点应互相对应，平移关系应为 $u=u'-2\Lambda^2$，相应奇点位于 $u'=\pm2\Lambda^2$。原文此处作 $u=u'-\Lambda^2/2$，与后文 $u'=\pm\Lambda^2/2$ 自洽，但与本节开头及图注的 $-4\Lambda^2$ 矛盾（v3 已改作 $u=u'-2\Lambda^2$）。按底本 v2 照排。]
这说明了场论非微扰分析中经常出现的一处微妙性。朴素地说，$u$ 定义为 $\vev{\tr\Phi^2/2}$。但是，复合算符需要通过仔细实施正则化和重整化来谨慎定义。由于超对称理论中两个手征算符之间几乎没有发散，对手征复合算符做这件事相对安全；不过，仍须注意两个依赖规范的手征算符之间的点分裂，例如它已知是 Konishi 反常的来源 \cite{Konishi:1983hf}。至少在微扰层面，可以采用全纯方案，它唯一确定了手征复合算符在微扰论所有阶的正则化和重整化。然而，在方案定义中仍存在非微扰歧义。在当前情形下，$u$ 与 $\Lambda^2$ 的质量量纲都为二，在已破缺的连续 R 对称性下的荷也都为 $2$，因此它们容易发生混合。在第 \ref{sec:su2nfpre} 小节猜测曲线时，我们没有计入未破缺的离散 R 对称性，因此曲线中出现的 $u$，与按定义构造成在离散 R 对称性下具有良好变换性质的 $u'$ 之间，存在差别。

我们得知，$u=0$ 与 $u=-\Lambda^2$，或等价地 $u'=\pm\Lambda^2/2$ 处的低能行为，由离散 R 对称性与味宇称的组合相互联系。让我们更详细地研究它们。已知 $u=0$ 处的单值变换由式 \eqref{nf_2_m+} 给出。设在 $u=0$ 附近，$a_D\sim cu$，其中 $c$ 为常数。施行一次 $S$ 变换，得到对偶耦合的跑动
\begin{equation}
\tau_D(E)\simeq+\frac{2}{2\pi i}\log\LambdaRG
\end{equation}
其中 $E\sim cu$ 确定能标。将其与纯 $\SU(2)$ 理论单极子点处的对偶耦合跑动 \eqref{dual_run} 比较。分子中的因子 $2$ 来自式 \eqref{nf_2_m+} 中 $M_+$ 的左下角矩阵元；更物理地说，它来自图 \ref{fig:collision_nf_2} 所示的两对分支点同时碰撞这一事实。一般而言，当 $\U(1)$ 规范理论与若干荷为 $q_i$ 的超多重态耦合时，跑动为
\begin{equation}
\tau\simeq+\frac{\sum_iq_i^2}{2\pi i}\log\LambdaRG\label{general_u1_running}
\end{equation}
因此，可以唯一地得出结论：存在两个荷为 $1$ 的超多重态。这一点也可以从高维观点看出：对于每一对碰撞的分支点，都有如图 \ref{fig:monopole_as_membrane} 那样的圆盘状膜。分支点真正碰撞时，它们变为无质量，从而提供两个带荷超多重态。

\subsection{$N_f=2$：模空间}\index{VacuumModuliSpace@真空模空间|hyperpage}\index{HiggsBranch@希格斯支|hyperpage}
在第 \ref{sec:hkq} 小节，我们研究过，含有两个荷为 $1$ 的超多重态的 $\U(1)$ 规范理论具有形如 $\bC^2/\bZ_2$ 的希格斯支。加上描述库仑支的 $u$ 平面，就可以将整个超对称真空模空间描绘为图 \ref{fig:nf_2_quantum_moduli}。注意，$u$ 平面上与希格斯支相接的两个奇点，会被离散 R 对称性与味宇称交换。
\begin{figure}[h]
\[
\inc{su2nf2-quantummoduli}
\]
\caption{$N_f=2$ 理论的量子模空间。\label{fig:nf_2_quantum_moduli}}
\end{figure}

将其与含有 $N_f=2$ 个味的 $\SU(2)$ 理论的经典模空间比较。库仑支仍由 $u=\tr\Phi^2/2$ 描述。当 $\Phi=0$ 时，可以进入希格斯支；我们在第 \ref{sec:hkq} 小节也研究过这个系统，看到它由 $\bC^2/\bZ_2\wedge\bC^2/\bZ_2$ 给出。可以将其描绘为图 \ref{fig:nf_2_classical_moduli}。
\begin{figure}[h]
\[
\inc{su2nf2-classicalmoduli}
\]
\caption{$N_f=2$ 理论的经典模空间。\label{fig:nf_2_classical_moduli}}
\end{figure}

在第 \ref{sec:hypergeneral} 小节，我们论证了规范动力学不能修正希格斯支的局部度规。这里可以看到，量子动力学仍然能够将两个 $\bC^2/\bZ_2$ 副本与 $u$ 平面相接的点分裂开；那一小节的论证不适用于度规奇异的点。

回忆一下，存在味对称性 $\SO(4)\simeq\SU(2)_A\times\SU(2)_B$，使得 $\SU(2)_{A,B}$ 分别作用于两个 $\bC^2/\bZ_2$ 副本。经过非微扰修正后，$\SU(2)_A$ 作用于 $u'=\Lambda^2/2$ 处的超多重态，$\SU(2)_B$ 则作用于 $u'=-\Lambda^2/2$ 处的超多重态。这与味宇称交换 $u'=\pm\Lambda^2/2$ 的作用相符，参见式 \eqref{flavor_parity}。

在第 \ref{sec:monopole} 小节，我们了解到，这类理论中的单极子按 $\SO(2N_f)$ 味对称性的旋量表示变换。这里，$\SO(4)$ 的旋量就是 $\SU(2)_A$ 或 $\SU(2)_B$ 的基本二重态，它们确实由味宇称交换。这与迄今的结果一致。
\begin{figure}[h]
\[
\inc{su2nf2-moduli-equalmass}
\]
\caption{当 $\mu_1=\mu_2$ 时，$N_f=2$ 理论的模空间。\label{fig:nf_2_moduli}}
\end{figure}
结束本小节之前，讨论将 $\mu=\mu_1=\mu_2$ 取为一个很小但非零的值时会发生什么。这将味对称性 $\SO(4)=\SU(2)_A\times\SU(2)_B$ 破缺为，例如 $\SU(2)_A$。相应地，直接分析曲线的判别式，可以检验原来同处于 $u=-\Lambda^2$ 的两个奇点会分裂为两个。$u$ 平面上仍有一个点，两个奇点在那里保持重合，其周围的局部单值变换仍为 $M_+$，没有改变。那里有一个形如 $\bC^2/\bZ_2$ 的希格斯支。所得结构如图 \ref{fig:nf_2_moduli} 所示。当连续增大 $\mu$ 时，最终更适合将这一情形视为图 \ref{fig:nf_2_u_plane_generic} 中 $\mu=\mu_1=\mu_2$ 的特殊情况。即在能标 $\mu$ 处，规范耦合仍然很弱，经典拉格朗日量分析有效。超势为
\begin{equation}
(Q_i^1,Q_i^2)\begin{pmatrix}\mu+a&0\\0&-a+\mu\end{pmatrix}
\begin{pmatrix}\tilde Q^i_1,\tilde Q^i_2\end{pmatrix}\label{supernf}
\end{equation}
因此，当 $a=\mu$ 时，$i=1,2$ 的分量 $(Q_i^2,\tilde Q^i_2)$ 保持无质量。规范群从 $\SU(2)$ 破缺为 $\U(1)$，并有两个荷为 $1$ 的超多重态，产生形如 $\bC^2/\bZ_2$ 的希格斯支。

\subsection{$N_f=3$}\label{sec:nf3}\index{SU2ThreeFlavors@含三个味的 $\SU(2)$ 理论|hyperpage}
含有 $N_f=3$ 个味的 $\SU(2)$ 理论的曲线，已在式 \eqref{nf_3_curve_pre} 中猜出：
\begin{equation}
\Sigma:\qquad\frac{(x-\tilde\mu_1)(x-\tilde\mu_2)}z+{2\Lambda(x-\tilde\mu_3)}z=x^2-u.\label{nf3curve!}
\end{equation}
可以看到，无论 $u$ 如何取值，$\lambda$ 都在 $z=0,1,\infty$ 处发散，并有四个随 $u$ 改变而移动的分支点，见图 \ref{fig:nf_3_curve}。为什么给质量参数加上波浪号，很快就会清楚。
\begin{figure}[h]
\[
\inc{su2nf3-curveX}
\]
\caption{$N_f=3$ 理论的曲线。\label{fig:nf_3_curve}}
\end{figure}

检验 $|u|\gg|\mu_i|^2,|\Lambda|^2$ 时的行为。两个分支点位于 $z\sim O(1)$，另一个分支点位于 $\sqrt{u}/\Lambda$。现在将 $A$ 闭链放在 $|z|=c$ 附近，其中 $1\ll c\ll\sqrt{u}/\Lambda$。积分与此前一样给出
\begin{equation}
a=\frac{1}{2\pi i}\oint_Ax\frac{dz}z\simeq\sqrt{u}.
\end{equation}
$B$ 闭链积分可以近似为
\begin{equation}
a_D\sim\frac{2}{2\pi i}\int^1_{\sqrt{u}/\Lambda}a\frac{dz}z\simeq-\frac{2}{2\pi i}a\log\frac{a}{\Lambda}.
\end{equation}
由此得到
\begin{equation}
\tau(a)=\frac{\partial a_D}{\partial a}=-\frac{2}{2\pi i}a\log\frac{a}{\Lambda},
\end{equation}

% [Tutor Note: 由上一式 $a_D\simeq-\frac{2}{2\pi i}a\log\frac{a}{\Lambda}$ 对 $a$ 求导，应为 $\tau(a)\simeq-\frac{2}{2\pi i}\bigl(\log\frac{a}{\Lambda}+1\bigr)$；原文（v2、v3 同）印作 $-\frac{2}{2\pi i}a\log\frac{a}{\Lambda}$，多出因子 $a$，疑为笔误。其渐近的对数系数仍与单圈计算 \eqref{general_1_loop} 相符。照原文排。]
从而重现了场论单圈计算 \eqref{general_1_loop}。
\begin{figure}[h]
\[
\inc{su2nf3-uplane-equalmasses}
\]
\caption{质量相等时，$N_f=3$ 理论的 $u$ 平面。\label{fig:nf_3_u_plane_generic}}
\end{figure}
当 $\mu_1=\mu_2=\mu_3=\mu$ 且 $|\mu|\gg|\Lambda|$ 时，能标 $\mu$ 处的耦合仍然很小，利用超势 \eqref{supernf} 的经典分析基本有效。预期在 $a\simeq\mu$ 附近，即 $u\simeq\mu^2$ 时，规范群 $\SU(2)$ 破缺为 $\U(1)$，并有三个荷为 $1$ 的超多重态。$u$ 平面上的这一点算作三个奇点，因为当 $\mu_{1,2,3}$ 略有不同时，它们应当位于三个略有区别的点 $u\simeq\mu_i^2$。当 $|u|\ll|\mu|^2$ 时，理论可以有效地用纯 $\SU(2)$ 规范理论描述，它具有单极子点和双荷子点。因此，总共预期 $u$ 平面上有五个奇点，见图 \ref{fig:nf_3_u_plane_generic}。

我们希望研究无质量情形 $\mu=0$。这里，不能直接在曲线 \eqref{nf3curve!} 中令 $\tilde\mu_i=0$。\footnote{作者感谢 Kazuya Yonekura 指出这一点。} 我们已经看到，当 $N_f=2$ 时，真空期望值 $u$ 可以与单瞬子因子 $\Lambda^2$ 混合。这里 $N_f=3$，单瞬子因子为 $\Lambda$，它可以与任何中性、手征、量纲为 $1$ 的算符混合。曲线只显现出 $\U(3)$ 味对称性。对应 $\U(1)$ 味对称性的质量参数是中性的、手征的，且量纲为 $1$。因此可以有形如
\begin{equation}
\tilde\mu_i=\mu_i+c\Lambda
\end{equation}
的混合，其中 $c$ 为常数。这里，我们规定不带波浪号的质量参数 $\mu_i$ 在 $\SO(2N_f)=\SO(6)$ 味对称性的 Weyl 对称性 $\mu_i\to\pm\mu_i$ 下线性变换。

为确定 $c$，令
\begin{equation}
(\tilde\mu_1,\tilde\mu_2,\tilde\mu_3)=(-\mu+c\Lambda,\mu+c\Lambda,\mu+c\Lambda)
\end{equation}
并研究 $u$ 平面上的奇点。这只是对具有 $\SU(3)$ 味对称性的质量选择所作的 $\SO(6)$ 味 Weyl 变换，因此 $u$ 平面上的五个奇点中，仍应有三个像图 \ref{fig:nf_3_u_plane_generic} 中那样碰撞。显式计算表明，这只有在 $c=1$ 时才会发生。
\begin{figure}[h]
\[
\inc{su2nf3-uplane-masslessX}
\]
\caption{无质量 $N_f=3$ 理论的 $u$ 平面。\label{fig:nf_3_uplane_massless}}
\end{figure}

终于可以令 $\mu=0$。现在曲线为
\begin{equation}
\frac{(x-\Lambda)^2}z+2\Lambda(x-\Lambda)z=x^2-u.
\end{equation}
$x(z)$ 在 $z=1$ 处有一个与 $u$ 无关的分支点。另外两个分支点随 $z$ 移动，位于方程
\begin{equation}
\Lambda^2z^2-\Lambda^2z+u-\Lambda^2=0.
\end{equation}
的解处。当 $u=\Lambda^2$ 或 $u=(5/4)\Lambda^2$ 时，分支点发生碰撞：
\begin{itemize}
\item 当 $u=(5/4)\Lambda^2$ 时，两个依赖 $u$ 的分支点在 $z=1/2$ 处相遇。那里的局部物理只是含有一个带荷超多重态的 $\U(1)$ 规范理论。
\item 当 $u=\Lambda^2$ 时，一个分支点移动到 $z=0$，另一个分支点与 $z=1$ 处不依赖 $u$ 的分支点碰撞。根据一般分析，$u$ 平面上有五个奇点，因此这一点应算作四个相碰的奇点，见图 \ref{fig:nf_3_uplane_massless}。
\end{itemize}
无穷远处的单值变换为
\begin{equation}
M_\infty=\begin{pmatrix}-1&1\\0&-1\end{pmatrix}.
\end{equation}
将绕 $u=\Lambda^2$、$u=(5/4)\Lambda^2$ 的单值变换记为 $M_+$、$M_-$，有
\begin{equation}
M_\infty=M_+M_-,\qquad M_+=\begin{pmatrix}1&0\\-4&1\end{pmatrix},\quad
M_-=\begin{pmatrix}-1&1\\-4&3\end{pmatrix}\sim\begin{pmatrix}1&-1\\0&1\end{pmatrix}.
\end{equation}

在 $u=\Lambda^2$ 处转到 S 对偶描述，得到对偶耦合的跑动
\begin{equation}
\tau(\LambdaRG)=+\frac{4}{2\pi i}\log\LambdaRG
\end{equation}
其中能标由 $\LambdaRG\sim(u-\Lambda^2)$ 确定。与式 \eqref{general_u1_running} 比较，可以猜测低能物理是一个 $\U(1)$ 规范理论，并且要么 (i) 只与一个荷为 $2$ 的超多重态耦合，要么 (ii) 与四个荷为 $1$ 的超多重态耦合。

回忆一下，经典理论具有希格斯支。选择 (i) 在 $u=\Lambda^2$ 处没有希格斯支，在 $u=(5/4)\Lambda^2$ 处也没有。量子修正应当保持希格斯支，因此这种选择被排除。

选择 (ii) 在 $u=\Lambda^2$ 处确实具有希格斯支。四个荷为 $1$ 的超多重态与 $\U(1)$ 规范多重态耦合。因而，希格斯支的复维数为 $2\cdot4-2\cdot1=6$。旋转这四个超多重态的 $\SU(4)$ 味对称性作用在其上。

在经典层面，有三个 $\SU(2)$ 二重态超多重态，因此希格斯支的复维数为 $4\cdot3-2\cdot3=6$，与上述计算一致。回忆一下，三个 $\SU(2)$ 二重态超多重态相当于六个 $\SU(2)$ 二重态半超多重态，具有 $\SO(6)$ 味对称性。由于 $\SO(6)\simeq\SU(4)$，希格斯支的对称性也相符。还应记住，这一理论中的单极子按 $\SO(2N_f)$ 味对称性的旋量表示变换。在当前情形中，$\SO(6)$ 的旋量是 $\SU(4)$ 的四维基本表示。这也与我们的选择一致：在 $u=\Lambda^2$ 处，存在四个与对偶 $\U(1)$ 电耦合的带荷超多重态。更详细的分析可以验证，这些希格斯支作为超凯勒流形也是一致的。
